\section{Anatomical basis for the expected-organ table}
\label{app:anatomy}

The correction described in \S\ref{sec:results:lesion} changes a headline
number, so the reasoning behind it is reproduced here rather than only cited. The
full source list, with quoted passages and the questions put to each source, is
in \code{docs/references/anatomy-citations.md} in the repository.

\paragraph{The question, asked before the change was made.}
Should the \code{EXPECTED} entries for T11 and T12 include the lungs? Put
anatomically: at typical inspiration, does the posterior lung base, via the
costophrenic recess, reach the T11 vertebral level, and the T12 level? Does this
differ by side, sex, or respiratory phase? The research was completed and the
answer fixed before the expected-organ table was edited and before the
validation was re-scored. This ordering matters, because the alternative
procedure, changing the table and keeping whichever version improves the metric,
would make the evaluation circular.

\paragraph{Primary quantitative source.}
Oh et al.\ \cite{oh2009retrodiaphragmatic} retrospectively reviewed 540
consecutive posteroanterior chest radiographs from general health screening, of
which 282 met inclusion criteria (clear visualization of the inferior boundary
of the posterior costophrenic sulcus behind the right hemidiaphragm, no
diaphragm-affecting abnormality; 155 men, 127 women, mean age
$40.7 \pm 8.4$ years). They measured the depth of the posterior costophrenic
sulcus below the diaphragmatic dome and scored its level against the spine.

The mean right posterior costophrenic sulcus depth was $29.2 \pm 15.6$\,mm, at
an average spinal level of $13.5 \pm 0.6$, which is the lower half of the L1
vertebral body: a full vertebral level below T12. Depth increased with height,
BMI and total diaphragmatic height, indicating normal population variation
rather than a rare tail. The authors' own conclusion is that a standard
posteroanterior chest radiograph must extend more than 3\,cm below the lateral
costophrenic sulcus, or to the L1 spine level, to capture the retrodiaphragmatic
lung fully.

Both T11 and T12 therefore sit at or above the population-average posterior
lung and pleural extent. A lung detected at either level is consistent with
normal anatomy and is not evidence of a phantom detection.

\paragraph{Triangulating sources.}
Three further references were consulted and agree on direction. A textbook-derived
thorax chapter states that lung tissue does not normally extend much below T10 in
quiet respiration, while the pleural cavity's inferior extent reaches
approximately T12 through the costodiaphragmatic recess, or lower along the
twelfth rib in the midscapular line, and flags these levels as approximate. A
descriptive anatomy reference for the costodiaphragmatic recess notes that during
deep inspiration the lungs descend up to two intercostal spaces into the recess.
A clinical teaching reference places the lower border of the lung at the
vertebral column at the upper half of T11 in moderate respiration, before any
inspiratory excursion is added.

The distinction that resolves the apparent disagreement between the T10 and
T11--L1 figures is that lung parenchyma in quiet breathing and the aerated lung
at inspiration are not the same boundary. The expected-organ table scores whether
a predicted lung mask could plausibly extend to a level, which tracks the
inspiratory and recess extent, not the resting position.

\paragraph{Points of convergence.}
Every source giving a specific level places the posterior lung base at T11 or
below: T11 in moderate respiration, approximately T12 for the recess, and the
lower half of L1 as a quantitative radiographic average. No source reported a
systematic left--right asymmetry large enough to justify treating the two lungs
differently. The Oh et al.\ study measured the right side specifically because
the right costophrenic sulcus is the one clearly visible below the liver and
diaphragm on a frontal film, which is a measurement-visibility constraint and
not a claim about the left side.

\paragraph{Decision, and what it changed.}
Both T11 and T12 include \org{lung-left} and \org{lung-right}. The weakest
single-source claim already places T11 within range, and T12 sits comfortably
below every conservative estimate. No source was found that would support
treating a lung detection at either level as anatomically impossible.

Applying the correction moved mean rank from 1.94 to 1.78, reduced anatomically
impossible top-1 detections from 38 to 34, and moved chance from 3.01 to 2.91,
since chance depends on expected-set size. It also moved T11 from below chance to
above it. As noted in \S\ref{sec:discussion}, the fact that one such error existed
is itself evidence that the table may contain others, particularly in this same
thoracolumbar region where the anatomy is genuinely variable.
